\setcounter{section}{7}
\section{Mapeos Conformes}
 
\begin{definition}
    \(U\ab, V \ab \subset \C\) son \emph{conformemente equivalentes}, se \(\exists f : U \to V \)  holomorfa e biyectiva (\emph{mapeo conforme}), escribimos \(f: U \simeq V \).  
\end{definition}

\begin{proposition}
    \textbf{1.1.} Se \(f \in \ho(U)\) es inyectiva, entonces \(f'(z) \neq 0, \ \forall z \in U\). En particular, la inversa \(f^{-1}: \operatorname{Im}f \to U\) es holomorfa, esto es, tambien es un mapeo conforme. 
\end{proposition}

\demo{Se \(f'(z_0) = 0 \), entonces \(f(z) -f(z_0) = a_k(z-z_0)^k + o(z^{k+1}) :  k\geq 2, \ a_k \neq 0 \), para \(0 < |w| \ll \epsilon \), tenemos \(f(z) - f(z_0) -w = [F(z) - w ] + o(z^{k+1})\), con al menos \(2\) soluciones distintas, \(\blacksquare\) \hyperlink{thm343}{4.3. (\emph{Rouché})}.\Lightning \ Ahora, siendo \(g:= f^{-1}: \operatorname{Im}f \to U\), 
\[\frac{g(w) -g(w_0)}{w-w_0} = \left(\frac{w -w_0}{g(w) -g(w_0)}\right)^{-1} = \left(\frac{f(z) - f(z_0)}{z-z_0}\right)^{-1} = \frac{1}{f'(z_0)}. \] }

\begin{definition}
    \centering \(\mathbb{H}:= \{z \in \C : \operatorname{Im}(z)>0\}\). \ \textcolor{gray}{\(\leftarrow\) \emph{Semi-plano superior}}
\end{definition}

\newcommand{\Ha}{\mathbb{H}}
\begin{theorem}
    \textbf{1.2.} \(\Ha \simeq \D\) vía \(\displaystyle F:  z \mapsto \frac{i-z}{i+z}\) e \(\displaystyle F^{-1}: w \mapsto i\  \frac{1-w}{1 + w} \). \comentario{Mogollas, faz conta}
\end{theorem}

\begin{definition}
    \(\displaystyle f : z \mapsto \frac{az + b}{cz + d} : a,b,c,d \in \C \ \) es una \emph{tranformación de Möbius}.
\end{definition}

\begin{note}
    En \cite[pags. 208-213]{stein} hay varios ejemplos de algunas de estas equivalencias. 
\end{note}

\subsection*{Lema de Schwarz e Automorfismos de \(\D\simeq \Ha\)}

\begin{lemma}
    \textbf{2.1. (\emph{Schwarz})} Sean \(f: \D \to \D \) tal que \(f(0) = 0 \). Entonces, 
    \vspace{-0.3cm}
    \begin{enumerate}[label = \roman*., left = 0.2cm ]
        \item \(|f(z)| \leq |z|, \ \forall z \in \D \). \hfill ii. \ Si \(\exists z_0 \in \D^\times : |f(z_0)| = |z_0| \Rightarrow f\) es \emph{rotación}. 
        \item[iii.] \(|f'(0)| \leq 1\), si valiera la igualdad entonces \(f\) es rotación.
    \end{enumerate}
\end{lemma}

\demo{(i) \ \(\nicefrac{|f(z)|}{|z|} \leq \nicefrac{1}{r}\) si \(|z| \leq r < 1\), singularidad removible + \(\blacksquare\) \hyperlink{thm245}{4.5. (\emph{PMM})}; (ii) \ \(\nicefrac{f(z)}{z}\) alcanza un máximo en \(\D\) luego \(f(z) = cz\) e \(|c| = 1 \ \leftrightharpoons \ f(z) = e^{i\theta} : \theta \in [0,2\pi]\); (iii) \ \(g(0) := \lim_{z\to 0}\nicefrac{f(z)}{z} = \lim_{z\to 0} \nicefrac{f(z) - f(0)}{z} = f'(0) = 1\), }

\begin{definition}
    Sea \(\Omega\ab\subset \C\), denotamos el \emph{grupo de automorfismos} de \(\Omega\) en sí mismo como \(\operatorname{Aut}(\Omega):= \{f: \Omega \to \Omega \ \mid \ f \text{ es conforme}\}\). También, dado \(\alpha \in \D\) definimos, 
    \[\frac{\alpha - z}{1 - \overline{\alpha} z} =: \psi_\alpha(z) \in \operatorname{Aut}(\D)\ \  \leftarrow  \text{\emph{Operadores de Blaschke}}. \]  
\end{definition}

\begin{theorem}
    \textbf{2.2.} Si \(f \in \operatorname{Aut}(\D)\), entonces \(\exists \theta \in \R\) e \(\exists \alpha \in \D  : f(z) = e^{i\theta}\cdot \psi_\alpha(z)\).
\end{theorem}

\demo{\(\exists ! \alpha \in \D : f (\alpha) = 0\), haga \(g(z):= (f \circ \psi_\alpha)(z)\). Note que \(g, g^{-1} \in \operatorname{Aut}(\D)\), \(g(0) = 0\), por el {\scriptsize \(\square\)} \hyperlink{lemma821}{2.1. (\emph{Schwarz})} \(|g(z)| = |z|, \ \forall z \in \D\), luego \(g(z) = e^{i\theta}z\).}% = e^{i\theta} \cdot \psi_\alpha(z)\).}

\begin{corollary}
    \textbf{2.3.} Las rotaciones son los únicos automorfismos en \(\D\) que fijan el origen. %os automorfismos en \(\D\) que fijan el origen son rotaciones.
\end{corollary}

\begin{note}
    \(\Gamma: \operatorname{Aut}(\D) \cong \operatorname{Aut}(\Ha)\) vía conjugación, \(\Gamma : \varphi \mapsto F^{-1} \circ \varphi \circ F\), siendo \(F\) el mapa conforme del \(\blacksquare\) \hyperlink{thm812}{1.2.}. Las transformaciones de Möbius están en \(\operatorname{Aut}(\Ha)\).%\(f_M\in \operatorname{Aut}(\Ha)\) es una transformación de Möbius,
    \[f_M(z) = \frac{az+ b}{cz+ d} \ \ \Big| \ \ M:= \begin{bmatrix}
    a & b \\ c & d   
    \end{bmatrix} \in \operatorname{SL}_2(\R) \leq \operatorname{GL}_2(\R).\]
\end{note}

\begin{theorem}
    \textbf{2.4.} \(\varphi \in \operatorname{Aut}(\Ha) \ \leftrightharpoons \ \exists M \in \operatorname{SL}_2(\R)  : \varphi = f_M\). \comentario{Véase \cite[pag. 222]{stein}}
\end{theorem}

\subsection*{Teorema de la Aplicación de Riemann}

\begin{theorem}
    \textbf{3.1. (\emph{Riemann Mapping Theorem})} Sean \(\Omega\ab \subsetneq \C\) simplemente conexo e \(z_0 \in \Omega\), entonces \(\exists ! F : \Omega \simeq \D\) tal que \(F(z_0)=0\) e \(F'(z_0)>0\).
\end{theorem}

